\documentclass[11pt,a4paper]{article}

\usepackage[utf8]{inputenc}
\usepackage[T1]{fontenc}
\usepackage[french]{babel}
\usepackage{lmodern}
\usepackage{geometry}
\geometry{margin=2.2cm}
\usepackage{amsmath,amssymb}
\usepackage{booktabs}
\usepackage{hyperref}
\usepackage{listings}
\usepackage{xcolor}
\usepackage{enumitem}
\usepackage{caption}
\usepackage{fancyhdr}

\hypersetup{
  colorlinks=true,
  linkcolor=blue!50!black,
  urlcolor=blue!50!black
}

\lstset{
  basicstyle=\ttfamily\footnotesize,
  breaklines=true,
  frame=single,
  columns=fullflexible,
  showstringspaces=false,
  keepspaces=true,
  backgroundcolor=\color{gray!5},
  extendedchars=true,
  literate=
    {é}{{\'e}}1 {è}{{\`e}}1 {ê}{{\^e}}1 {ë}{{\"e}}1
    {à}{{\`a}}1 {â}{{\^a}}1
    {ù}{{\`u}}1 {û}{{\^u}}1
    {ô}{{\^o}}1 {ö}{{\"o}}1
    {î}{{\^i}}1 {ï}{{\"i}}1
    {ç}{{\c{c}}}1
    {É}{{\'E}}1 {È}{{\`E}}1 {À}{{\`A}}1 {Ç}{{\c{C}}}1
    {–}{{-}}1 {—}{{--}}1
}

\pagestyle{fancy}
\fancyhf{}
\fancyhead[L]{\small PredLDA --- Int\'egration \& livraison}
\fancyhead[R]{\small \'Equipe G}
\fancyfoot[C]{\thepage}
\renewcommand{\headrulewidth}{0.4pt}

\title{%
  PredLDA\\[0.35em]
  \large Mise en production, int\'egration DevOps\\ et livraison du package R\\[0.4em]
  \normalsize Rapport d'int\'egration --- M2 SISE Programmation R 2026--2027}
\author{Thomas Fonferrier\\[0.2em]
{\small Avec Ryan Ducret, Aya Saadi, Xilian Jing --- \'Equipe G}}
\date{\today}

\begin{document}
\maketitle

\begin{abstract}
Ce document d\'ecrit \textbf{tout le travail d'int\'egration} r\'ealis\'e pour
mettre en production et livrer le package R \texttt{PredLDA} : architecture du
d\'ep\^ot GitHub, m\'etadonn\'ees du package, documentation anglaise aux normes~R,
tutoriel reproductible, jeux d'exemple, application Shiny embarqu\'ee,
int\'egration continue, artefacts de livraison (\texttt{.tar.gz}, rapports) et
proc\'edure de validation d'installation. Il se concentre sur la
\emph{cha\^ine de livraison}, compl\'ementaire au rapport m\'ethodologique LDA.
\end{abstract}

\tableofcontents
\vspace{1em}

\section{Contexte et objectifs de l'int\'egration}

Le cahier des charges impose un package :
\begin{itemize}[leftmargin=*]
  \item installable \textbf{directement depuis GitHub} ;
  \item document\'e en anglais aux normes R (\texttt{?LDA} comparable \`a \texttt{?glm}) ;
  \item accompagn\'e d'un \textbf{tutoriel reproductible} et de jeux d'exemple ;
  \item livr\'e avec une application Shiny, un rapport LaTeX et un plan~B hors
        ligne (archive \texttt{.tar.gz}).
\end{itemize}

Le r\^ole d'int\'egration (DevOps l\'eger) consiste \`a transformer un code
algorithmique et une application Shiny en un \textbf{produit R livrable} :
structure standard, m\'etadonn\'ees valides, documentation, tests, CI et
proc\'edure de v\'erification \og comme le destinataire final \fg{}.

\begin{table}[ht]
\centering
\small
\begin{tabular}{@{}lp{9.2cm}@{}}
\toprule
Notion & \'Equivalent concret ici \\
\midrule
Build & \texttt{R CMD build} / \texttt{devtools::build()} \\
Contr\^ole qualit\'e & \texttt{R CMD check} / \texttt{devtools::check()} \\
Distribution & GitHub + \texttt{remotes::install\_github} \\
CI & GitHub Actions (\texttt{R-CMD-check}) \\
Artefact plan~B & \texttt{PredLDA\_0.1.0.tar.gz} \\
Documentation runtime & pages Rd + vignette + README \\
\bottomrule
\end{tabular}
\caption{Correspondance DevOps $\leftrightarrow$ \'ecosyst\`eme package R.}
\end{table}

\section{Architecture du d\'ep\^ot et du package}

\subsection{Organisation GitHub}

Le d\'ep\^ot
\url{https://github.com/thomas-fonferrier/projet-lda-r}
h\'eberge le package dans le sous-dossier \texttt{PredLDA/} (et non \`a la
racine). Cette organisation impose le param\`etre
\texttt{subdir = "PredLDA"} lors de l'installation GitHub.

\`A la racine du d\'ep\^ot ont \'et\'e ajout\'es :
\begin{itemize}[leftmargin=*]
  \item \texttt{README.md} --- point d'entr\'ee (install, quick start, liens) ;
  \item \texttt{.gitignore} --- exclusion des artefacts locaux ;
  \item \texttt{.github/workflows/R-CMD-check.yml} --- int\'egration continue.
\end{itemize}

\subsection{Structure du package}

\begin{lstlisting}
PredLDA/
|-- DESCRIPTION          # identite, auteurs, dependances
|-- NAMESPACE            # exports (genere par roxygen2)
|-- LICENSE              # MIT
|-- README.md            # tutoriel EN reproductible
|-- R/
|   |-- lda_r6.R         # classe publique LDA + doc roxygen
|   +-- stats_utils.R    # briques statistiques internes
|-- man/                 # aide ?LDA (Rd)
|-- tests/               # tests unitaires / reference MASS
|-- vignettes/           # vignette predlda-intro
|-- inst/
|   |-- extdata/         # jeux d'exemple du tutoriel
|   +-- shiny/           # application Shiny + examples/
|-- data-raw/            # scripts de sync des exemples
+-- report/              # sources LaTeX (exclus du build R)
\end{lstlisting}

Points d'int\'egration importants :
\begin{itemize}[leftmargin=*]
  \item \texttt{report/} est list\'e dans \texttt{.Rbuildignore} : le rapport
        n'est pas embarqu\'e dans le package install\'e ;
  \item les placeholders vides \texttt{ui.R}/\texttt{server.R}/\texttt{global.R}
        ont \'et\'e supprim\'es : l'app est enti\`erement dans \texttt{app.R} ;
  \item \texttt{inst/extdata/} expose les m\^emes fichiers d'exemple que Shiny
        via \texttt{system.file("extdata", ...)}.
\end{itemize}

\section{M\'etadonn\'ees DESCRIPTION et correction Authors@R}

\subsection{Contenu livr\'e}

Le fichier \texttt{DESCRIPTION} d\'eclare le titre, la version \texttt{0.1.0},
la licence MIT, les quatre auteurs de l'\'equipe~G, \texttt{Imports: R6}, les
\texttt{Suggests} (MASS, shiny, readxl, knitr, rmarkdown),
\texttt{VignetteBuilder: knitr}, ainsi que \texttt{URL}/\texttt{BugReports}.

\subsection{Incident de production corrig\'e}

Une premi\`ere installation GitHub a \'echou\'e avec le message :
\begin{lstlisting}
Authors@R field gives no person with maintainer role,
valid email address and non-empty name.
\end{lstlisting}
Cause : syntaxe R invalide.
\begin{lstlisting}
# Incorrect : "cre" n'est PAS un role
role = "aut", "cre"

# Correct
role = c("aut", "cre")
\end{lstlisting}
Sans le r\^ole \texttt{cre} (maintainer), \texttt{R CMD build} refuse le package.
Correction pouss\'ee sur GitHub : Ryan Ducret est \texttt{[aut, cre]}.

\textbf{Le\c{c}on DevOps :} valider \texttt{Authors@R} avant tout push destin\'e
\`a \texttt{install\_github}.

\section{Documentation anglaise aux normes R}

La documentation est \'ecrite en roxygen2 dans \texttt{R/lda\_r6.R}, puis
r\'eg\'en\'er\'ee par :
\begin{lstlisting}
devtools::document("PredLDA")
\end{lstlisting}
Cela produit \texttt{man/LDA.Rd}, \texttt{man/PredLDA-package.Rd} et met \`a jour
le \texttt{NAMESPACE}.

La page \texttt{?LDA} documente d\'esormais :
\begin{itemize}[leftmargin=*]
  \item l'API (\texttt{new}, \texttt{fit}, \texttt{predict},
        \texttt{predict\_proba}, \texttt{score}, \texttt{print}, \texttt{summary}) ;
  \item les param\`etres et champs expos\'es ;
  \item une section \textbf{Reading the results} (une colonne = une classe,
        intercept en premi\`ere ligne, softmax, matching par nom, cas binaire
        $D(x)$) ;
  \item des exemples ex\'ecutables, y compris un fichier \texttt{extdata}.
\end{itemize}

\section{Tutoriel reproductible et jeux d'exemple}

Trois niveaux :
\begin{enumerate}[leftmargin=*]
  \item README racine du d\'ep\^ot --- installation GitHub et quick start ;
  \item \texttt{PredLDA/README.md} --- tutoriel EN complet ;
  \item vignette \texttt{predlda-intro} ---
        \texttt{vignette("predlda-intro", package = "PredLDA")}.
\end{enumerate}

Fichiers livr\'es dans \texttt{inst/extdata/} :
\texttt{iris\_apprentissage.tsv}, \texttt{iris\_apprentissage\_fr.csv},
\texttt{iris\_nouveaux.tsv}.

\begin{lstlisting}
ex <- system.file("extdata", "iris_nouveaux.tsv", package = "PredLDA")
newdata <- read.delim(ex, check.names = FALSE)
\end{lstlisting}

\section{Application Shiny embarqu\'ee}

L'application est livr\'ee sous \texttt{inst/shiny/} :
\begin{lstlisting}
shiny::runApp(system.file("shiny", package = "PredLDA"))
\end{lstlisting}
Le tutoriel Shiny EN d\'etaille un parcours reproductible (iris, split 80/20,
pr\'ediction individu/fichier). Les d\'ependances Shiny sont en
\texttt{Suggests}.

\section{Contr\^ole qualit\'e : tests, check et CI}

\subsection{Tests}

\begin{itemize}[leftmargin=*]
  \item \texttt{tests/test-LDA.R} --- API, validations, accord \texttt{MASS::lda} ;
  \item \texttt{tests/test-reference.R} --- formules textbook + hors-\'echantillon ;
  \item \texttt{inst/shiny/tests/test-server.R} --- logique serveur Shiny.
\end{itemize}

\subsection{R CMD check}

\begin{lstlisting}
devtools::check("PredLDA", args = c("--no-manual"))
# Status: OK
\end{lstlisting}

\subsection{GitHub Actions}

Le workflow \texttt{.github/workflows/R-CMD-check.yml} se d\'eclenche \`a chaque
\texttt{push}/\texttt{pull\_request} sur \texttt{main}, avec
\texttt{working-directory: PredLDA}, pour d\'etecter toute r\'egression avant
livraison.

\section{Cha\^ine de build et artefacts}

\subsection{Archive source (plan B)}

\begin{lstlisting}
devtools::build("PredLDA")
# -> PredLDA_0.1.0.tar.gz

install.packages("PredLDA_0.1.0.tar.gz",
                 repos = NULL, type = "source")
\end{lstlisting}

\subsection{Installation en ligne}

\begin{lstlisting}
remotes::install_github(
  "thomas-fonferrier/projet-lda-r",
  subdir = "PredLDA",
  upgrade = "never"
)
\end{lstlisting}

\subsection{Flux de mise en production}

\begin{enumerate}[leftmargin=*]
  \item Int\'egrer le code dans \texttt{PredLDA/} ;
  \item \texttt{devtools::document()} puis tests locaux ;
  \item \texttt{devtools::check()} jusqu'\`a \texttt{Status: OK} ;
  \item commit + push sur \texttt{main} ;
  \item v\'erifier GitHub Actions ;
  \item tester \texttt{install\_github} en session R propre ;
  \item reconstruire \texttt{.tar.gz} et PDF ;
  \item d\'eposer sur le drive + mail au responsable.
\end{enumerate}

\section{Proc\'edure de validation \og comme le professeur \fg{}}

Dans RStudio, sans d\'ependre du dossier local :

\begin{lstlisting}
try(remove.packages("PredLDA"), silent = TRUE)
remotes::install_github(
  "thomas-fonferrier/projet-lda-r",
  subdir = "PredLDA",
  upgrade = "never",
  force = TRUE
)
library(PredLDA)
?LDA
vignette("predlda-intro", package = "PredLDA")

model <- LDA$new()
model$fit(iris[, 1:4], iris$Species)
model$predict(iris[c(1, 51, 101), 1:4])

shiny::runApp(system.file("shiny", package = "PredLDA"))
\end{lstlisting}

Checklist : installation OK, \texttt{?LDA} lisible, vignette accessible,
\texttt{extdata} non vide, Shiny d\'emarre depuis le package install\'e.

\section{Inventaire des livrables d'int\'egration}

\begin{table}[ht]
\centering
\small
\begin{tabular}{@{}p{4.0cm}p{8.4cm}@{}}
\toprule
Livrable & Impl\'ementation r\'ealis\'ee \\
\midrule
Install GitHub & README + \texttt{subdir} + DESCRIPTION valide \\
Maintainer & Correction \texttt{Authors@R} (\texttt{c("aut","cre")}) \\
Aide EN & Roxygen enrichi + \texttt{man/*.Rd} \\
Tutoriel EN & README package + vignette \texttt{predlda-intro} \\
Jeux d'exemple & \texttt{inst/extdata/} + sync \texttt{data-raw/} \\
Shiny livrable & \texttt{inst/shiny/app.R}, via \texttt{system.file} \\
CI & GitHub Actions R-CMD-check \\
Plan B & \texttt{PredLDA\_0.1.0.tar.gz} build\'e et test\'e \\
Qualit\'e & \texttt{R CMD check} : Status OK \\
Rapports & \texttt{rapport.tex} + \texttt{rapport\_integration.tex} \\
\bottomrule
\end{tabular}
\caption{Synth\`ese des livrables d'int\'egration.}
\end{table}

\section{Limites restantes}

Hors p\'erim\`etre de cette int\'egration (\'equipe algo / suite) :
pr\'etraitement des variables cat\'egorielles, s\'election automatique de
variables, indicateurs de pertinence, enrichissement \'eventuel du rapport
m\'ethodologique (captures Shiny, r\'esultats d\'etaill\'es).

\section{Conclusion}

La mise en production de PredLDA repose sur une cha\^ine compl\`ete : d\'ep\^ot
GitHub structur\'e, m\'etadonn\'ees valides, documentation et tutoriel anglais,
exemples embarqu\'es, check R + CI, archive hors ligne et test d'installation en
session propre. Ces \'el\'ements permettent de \textbf{livrer} le package
conform\'ement au cahier des charges.

\vspace{1em}
\noindent\textbf{D\'ep\^ot :}
\url{https://github.com/thomas-fonferrier/projet-lda-r}

\appendix
\section{Commandes rapides}

\begin{lstlisting}
# Qualite locale
Rscript -e 'devtools::document("PredLDA")'
Rscript -e 'devtools::check("PredLDA", args="--no-manual")'
Rscript -e 'devtools::build("PredLDA", path=".")'

# Publication
git add -A
git commit -m "Release integration docs"
git push
\end{lstlisting}

\section{Trame du mail de rendu}

\begin{lstlisting}
A : ricco.rakotomalala@univ-lyon2.fr
Objet : [SISE - Prog. R - Equipe G] Fonferrier Ducret Saadi Jing

Bonjour,

Lien du projet PredLDA :
https://github.com/thomas-fonferrier/projet-lda-r

Livrables : package installable (subdir PredLDA), application Shiny,
tutoriel EN, archive PredLDA_0.1.0.tar.gz, rapports PDF/LaTeX.

Cordialement,
Equipe G
\end{lstlisting}

\end{document}
